@ID: OW16D1P1I
@BIDANG: Analisis Real

Jika

$$
S=\left\{\frac{n}{2m}+\frac{6m}{n}:n,m\in\mathbb N\right\},
$$

maka $\inf(S)=\ldots$ dan $\sup(S)=\ldots$.

@ID: OW16D1P2I
@BIDANG: Analisis Real

Misalkan $a_i>0$, $\forall i=1,2,\ldots,2016$. Jika

$$
(a_1a_2\cdots a_{2016})^{\frac1{2016}}=2,
$$

maka

$$
(1+a_1)(1+a_2)\cdots(1+a_{2016})\ge \ldots.
$$

@ID: OW16D1P3I
@BIDANG: Analisis Real

Jika barisan bilangan real $(x_n)$ memenuhi sifat

$$
\lim_{n\to\infty}(x_{2n}+x_{2n+1})=315
$$

dan

$$
\lim_{n\to\infty}(x_{2n}+x_{2n-1})=2016,
$$

maka

$$
\lim_{n\to\infty}\frac{x_{2n}}{x_{2n+1}}=\ldots.
$$

@ID: OW16D1P4I
@BIDANG: Analisis Real

Hitung

$$
\lim_{n\to\infty}\sum_{k=1}^{n}\frac{8n^2}{n^4+1}.
$$

@ID: OW16D1P5I
@BIDANG: Analisis Real

Untuk setiap $n\in\mathbb N$,

$$
f_n(x)=nx(1-x^2)^n,
$$

untuk setiap $x$, $0\le x\le 1$, dan

$$
a_n=\int_0^1 f_n(x)\,dx.
$$

Jika

$$
s_n=\sin(\pi a_n),
$$

untuk setiap $n\in\mathbb N$, maka

$$
\lim_{n\to\infty}s_n=\ldots.
$$

@ID: OW16D1P6I
@BIDANG: Analisis Real

Jika

$$
E=\{f\mid f:\mathbb R\to\mathbb R \text{ fungsi kontinu dengan } f(x)\in\mathbb Q,\ \forall x\in\mathbb R\},
$$

maka $E=\ldots$.

@ID: OW16D1P7I
@BIDANG: Analisis Real

Diberikan fungsi $f:\mathbb R\to\mathbb R$, dengan

$$
f(x)=4x,
$$

untuk bilangan rasional $x$, dan

$$
f(x)=x+6,
$$

untuk bilangan irasional $x$. Jika

$$
E=\{x\in\mathbb R:f \text{ kontinu di }x\},
$$

maka himpunan semua titik-limit $E$ adalah $\ldots$.

@ID: OW16D1P8I
@BIDANG: Analisis Real

Diketahui

$$
a<\frac{\pi}{2}.
$$

Jika $M<1$ dengan

$$
|\cos x-\cos y|\le M|x-y|,
$$

untuk setiap $x,y\in[0,a]$, maka

$$
M=\ldots.
$$
```

@ID: OW16D1P1U
@BIDANG: Analisis Real

Diketahui himpunan $E\subseteq\mathbb R$. Didefinisikan fungsi $f:\mathbb R\to\mathbb R$, dengan

$$
f(x)=\inf\{|x-y|:y\in E\},
$$

untuk setiap $x\in\mathbb R$. Buktikan bahwa fungsi $f$ kontinu pada $\mathbb R$.

@ID: OW16D1P2U
@BIDANG: Analisis Real

Diketahui fungsi $f:\mathbb R\to\mathbb R$ terdiferensial hingga tingkat $2$. Jika terdapat bilangan positif $A$ dan $B$, dengan

$$
|f(x)|\le A
$$

dan

$$
|f''(x)|\le B,
$$

untuk setiap $x\in\mathbb R$, buktikan bahwa

$$
|f'(0)|\le 2\sqrt{AB}.
$$

@ID: OW16D1P3U
@BIDANG: Analisis Real

Diketahui fungsi $f$ kontinu pada $[0,1]$. Didefinisikan $f_0=f$ pada $[0,1]$ dan untuk setiap $n\in\mathbb N$,

$$
f_{n+1}(x)=\int_0^x f_n(t)\,dt,
$$

untuk setiap $x\in[0,1]$.

Buktikan bahwa barisan fungsi $(f_n)$ konvergen seragam pada $[0,1]$ ke fungsi nol.
```

@ID: OW16D2P1I
@BIDANG: Kombinatorika

Sepotong kawat berukuran $1$ meter dipotong secara acak menjadi tiga bagian. Besarnya peluang ketiga bagian ini membentuk sebuah segitiga adalah $\ldots$.

@ID: OW16D2P2I
@BIDANG: Kombinatorika

Sebuah *palindrome* adalah sebuah barisan berhingga karakter sehingga dapat dibaca dengan cara yang sama baik dari kiri maupun kanan. SIKAPAKIS adalah sebuah *palindrome*. Banyaknya bilangan *palindrome* yang terdiri atas $7$ digit sedemikian sehingga tidak terdapat digit yang muncul lebih dari dua kali adalah $\ldots$.

@ID: OW16D2P3I
@BIDANG: Kombinatorika

Dari himpunan $26$ huruf

$$
A=\{a,b,c,\ldots,y,z\}
$$

dibentuk susunan enam huruf berbeda (susunan tak perlu bermakna) sedemikian sehingga huruf pertama dan huruf terakhir adalah huruf vokal, dan sisanya adalah huruf konsonan. Jika huruf $b$ selalu muncul pada susunan dan berdampingan dengan huruf $c$, maka banyaknya susunan yang mungkin adalah $\ldots$.

@ID: OW16D2P4I
@BIDANG: Kombinatorika

Banyaknya cara memfaktorkan bilangan $441.000$ menjadi dua faktor positif $m$ dan $n$ yang saling prima relatif adalah $\ldots$.

@ID: OW16D2P5I
@BIDANG: Kombinatorika

Banyaknya bilangan antara $1$ dan $500$ yang tidak habis dibagi oleh $3$, $4$, dan $6$ adalah $\ldots$.

@ID: OW16D2P6I
@BIDANG: Kombinatorika

Banyaknya graf sederhana berlabel atas $n$ titik yang memiliki sedikitnya dua sisi adalah $\ldots$.

@ID: OW16D2P7I
@BIDANG: Kombinatorika

Didefinisikan suatu fungsi rekursif, $\forall n\in\mathbb Z$, berlaku

$$
f(1)=1,\qquad
f(2)=5,
$$

dan

$$
f(n+1)=f(n)+2f(n-1),\qquad \forall n>2.
$$

Maka

$$
f(n)=\ldots.
$$

@ID: OW16D2P8I
@BIDANG: Kombinatorika

Dalam bentuk yang paling sederhana fungsi pembangkit biasa (*ordinary generating function*), $g(x)$, dari barisan

$$
(1,2,3,4,\ldots)
$$

adalah $\ldots$.

@ID: OW16D2P1U
@BIDANG: Kombinatorika

Perlihatkan bahwa untuk setiap himpunan yang terdiri dari $7$ bilangan bulat berbeda, maka terdapat dua bilangan $x$ dan $y$ pada himpunan tersebut sedemikian sehingga

$$
x+y
$$

atau

$$
x-y
$$

adalah kelipatan $10$.


